% discrete_measure_theory.tex
% Section 6 -- discrete measure theory on the bipartite octahedral
% lattice.  Counting measures, Haar-analogue measures invariant under
% the automorphism algebra, change-of-variables, discrete Radon-Nikodym.

\label{sec:discrete_measure_theory}

\placeholder{Section body -- write during Phase 3.  Develop the
measure theory needed to underwrite the function-space theory in
Section~\ref{sec:function_spaces} and the integration-by-parts
identity in Section~\ref{sec:discrete_diff_geometry}.}

\subsection{Counting Measure}
\label{subsec:counting_measure}

\placeholder{The basic counting measure $\mu_\#$ on $\Lambda_n$;
its restriction to colour classes $V_n^{\pm}$; sigma-algebra
generated by finite sets.}

\subsection{Haar-Analogue Measures}
\label{subsec:haar_analogue}

\placeholder{Measures invariant under the per-site automorphism
algebra of Section~\ref{sec:automorphism_algebra}; uniqueness up to
scaling; the analogue of Haar's theorem for the lattice.}

\subsection{Change-of-Variables}
\label{subsec:change_of_variables}

\placeholder{Change-of-variables formula for measure-preserving
maps and for the lattice-spacing rescaling $\epsilon \mapsto
\lambda \epsilon$.  Underwrites the limit-rescaling in
Section~\ref{sec:continuum_limits}.}

\subsection{Discrete Radon-Nikodym}
\label{subsec:radon_nikodym}

\placeholder{Discrete analogue of the Radon-Nikodym theorem: when a
measure $\nu$ on $\Lambda_n$ has a density with respect to
$\mu_\#$, and the explicit formula for the density.}
